\subsection{The Sequencing\textendash Whirling Limit Conjecture}

First stated in \texttt{dev/content/notes/2026/07/02/} at \texttt{metroplat.emsenn.net}
(``Sequencing\textendash Whirling Limit Conjecture''). Finite causal-diamond motion is
\texttt{sequencingPrec0} (\cref{thm:hole-diamond-kinematic-space}): acyclic, irreflexive endpoint
succession, $p \prec^\diamond q \iff p.\text{upper} \le q.\text{lower}$. The conjecture proposes
reading \emph{whirling} as what this acyclic succession becomes at a genuine limit stage.

\begin{theorem}[Sequencing\textendash Whirling Limit Conjecture]\label{thm:sequencing-whirling-limit-conjecture}\notready\uses{thm:hole-diamond-kinematic-space, thm:transfinite-closure-flow-no-blowup, thm:vortex-h1-obstruction}
\emph{(Open frontier --- stated, not formalized; no bridging Lean exists between the two sides.)}
Let $p_0 \prec^\diamond p_1 \prec^\diamond p_2 \prec^\diamond \cdots$ be an infinite sequencing
chain of causal-diamond points. At every finite stage the process is acyclic succession only. Two
already-proved modules describe the two possible behaviours at a limit stage $p_\omega$, but
nothing yet identifies them as the two branches of one dichotomy:
\begin{itemize}
  \item \textbf{Smooth limit.} \texttt{transfinite\_closure\_flow\_no\_blowup}
  (\cref{thm:transfinite-closure-flow-no-blowup}) proves that at the canonical limit
  $\omega = \texttt{Ord.limit}(\texttt{Ord.ofNat})$, the settled value is the directed join of the
  finite predecessor stages, with no non-Scott-continuous firing. The conjecture reads this as the
  \emph{no-whirl} case: $p_\omega$ remembers nothing beyond what the finite stages already
  determined.
  \item \textbf{Singular limit.} \texttt{VortexCohomology}'s discrete de Rham $H^1$
  (\cref{thm:vortex-h1-obstruction}) proves nonzero circulation is exactly what prevents a
  velocity cochain from being a gradient. The conjecture reads a nonzero-circulation obstruction at
  the limit stage as \emph{whirl}: $\text{Whirl}(p_\bullet) \iff \lnot\,\text{ContinuousAt}(p_\bullet,\omega)
  \iff [\text{circ}(p_\bullet)] \ne 0 \in H^1$.
\end{itemize}
The precise claim -- that these two already-formalized notions (Scott-continuity-at-$\omega$ from
one module, circulation-nonvanishing from an unrelated one) are literally dual descriptions of a
single limit-stage dichotomy for sequencing chains specifically -- has no supporting Lean
construction. Building it would require a genuine bridge object relating
\texttt{DiamondPoint}-valued transfinite chains to \texttt{VortexLoop}'s cochain carrier; neither
module currently mentions the other.
\end{theorem}
